\color{unidarkgreen}{\subsection{Поперечная дифференциальная защита реактора}}
\color{black}

\begin{enumerate}[label=\arabic{section}.\arabic{subsection}.\arabic*,  align=left, labelsep=4pt, leftmargin=0pt, itemindent=57pt]

\item
Поперечная дифференциальная токовая защита шунтирующего реактора (ПДЗ) предназначена
для защиты при возникновении витковых замыканий в сетевой обмотке, имеющей расщепление в каждой
из фаз со стороны нейтральных выводов.

\item
Защита реагирует на дифференциальные токи фаз А, В, С. Дифференциальный ток (Iдиф) в каждой
из фаз может быть получен путем программного вычисления векторной разности токов, измеренных в
обеих ветвях расщепленной обмотки соответствующей фазы реактора на стороне его нейтральных
выводов, или измерен, используя дифференциальный ТТ специального исполнения типа ДТФ либо
включением вторичных цепей на разность токов по схеме «восьмерка».

\item 
Защита выполнена с торможением, ток торможения (Iторм) может быть получен путем программного
вычисления векторной суммы токов, измеренных в обеих ветвях расщепленной обмотки
соответствующей фазы реактора на стороне его нейтральных выводов, либо измерен путем
электрического сложения токов нейтральных выводов.

\item
Выбор схемы подключения задается программной накладкой \textbf{«Подкл.»}. В случае контроля токов
в обеих ветвях расщепленной обмотки соответствующей фазы реактора на стороне его нейтральных
выводов (IНВ1, IНВ2), дифференциальные токи и токи торможения при этом рассчитываются как:

\begin{equation}
I_{\textsl{диф.х}} = \lvert K_{\textsl{ампл.1}} \cdot I_{\textsl{НВ1.х}} - K_{\textsl{ампл.2}} \cdot I_{\textsl{НВ2.х}} \rvert
\end{equation}

\begin{equation}
I_{\textsl{торм.х}} = \lvert K_{\textsl{ампл.1}} \cdot I_{\textsl{НВ1.х}} + K_{\textsl{ампл.2}} \cdot I_{\textsl{НВ2.х}} \rvert
\end{equation}

\begin{eqexpl}[5mm]
\eqexplSetItemAlign{r}  % выравнивание обозначений по левому краю(l)/правому краю(r)/центру(c)
\item{x} наименование фазы (A, B, C);
\end{eqexpl}

\item
В случае электрического подключения на разность (IДТФ) и сумму (IНВ) токов ветвей расщепленной
обмотки реактора дифференциальные токи и токи торможения, используемые функцией, определяются
по формуле:

\begin{equation}
I_{\textsl{диф.х}} = \lvert K_{\textsl{ампл.1}} \cdot I_{\textsl{ДТФ.х}} \rvert
\end{equation}

\begin{equation}
I_{\textsl{торм.х}} = \lvert K_{\textsl{ампл.2}} \cdot I_{\textsl{НВ.х}} \rvert
\end{equation}

\item
Защита имеет характеристику срабатывания, состоящую из двух участков: первый участок
горизонтальный (без торможения), второй участок наклонный (с торможением). Характеристика
срабатывания защиты описывается следующими параметрами:
\begin{itemize}
\item \textit{«Iср.диф»} -- минимальный дифференциальный ток срабатывания;
\item \textit{«Iнач.торм»} -- начальный ток торможения наклонного участка;
\item \textit{«Kторм»} -- коэффициент торможения наклонного участка.
\end{itemize}

\begin{figure}[H]
\centering
\includegraphics[width=0.4\textwidth,height=0.4\textheight,keepaspectratio]{\fbpath/161.0612 ПДЗт/img/char_srab.pdf}
\captionof{figure}{Характеристика срабатывания поперечной дифференциальной защиты реактора}\label{pict_pdzt:char_srab}
\end{figure}

\item
Значение срабатывания дифференциальной защиты с торможением на соответствующем участке
характеристики определяется по формулам:

\begin{itemize}
    \item участок 1: $I_{\textsl{диф}} \geq I_{\textsl{ср.диф}}$, при $I_{\textsl{торм}} \leq I_{\textsl{нач.торм}}$
    \item участок 2: $I_{\textsl{диф}} \geq I_{\textsl{ср.диф}} + K_{\textsl{торм}} \cdot (I_{\textsl{торм}} - I_{\textsl{нач.торм}})$, при $I_{\textsl{торм}} > I_{\textsl{нач.торм}}$
\end{itemize}

\begin{figure}[H]
\centering
\includegraphics[width=0.7\textwidth,height=0.7\textheight,keepaspectratio]{\fbpath/161.0612 ПДЗт/img/PDZt.pdf}
\captionof{figure}{Функциональный блок контроля ПДЗт}\label{pict_pdzt:PDZt}
\end{figure}


\small
\begin{longtable}{|>{\centering\arraybackslash}m{5.3cm}|>{\centering\arraybackslash}m{4.2cm}|>{\centering\arraybackslash}m{1cm}|>{\centering\arraybackslash}m{3.3cm}|>{\centering\arraybackslash}m{1.8cm}|}
% Первая страница: полный заголовок
\caption{Перечень уставочных параметров ПДЗт\hfill\vspace{-0.5\baselineskip}}\label{dtz:tbl2}\\
\hline
\rowcolor{gray!30}
Наименование уставки & Диапазон & Шаг & Значение по умолчанию & Единица измерения \\
\hline
\endfirsthead

% Промежуточные страницы: «Продолжение таблицы»
\caption*{\hspace{3pt}\emph{Продолжение таблицы \ref{dtz:tbl2}}\hfill\ }\vspace{-0.5\baselineskip} \\
\hline
\rowcolor{gray!30}
Наименование уставки & Диапазон & Шаг & Значение по умолчанию & Единица измерения \\
\hline
\endhead

%===m>ПДЗт|-
%===m
\end{longtable}
\normalsize

\end{enumerate}